% page 100 du fac-simile (image p-099 du scan)
% bloc 154.9 x 224.7 mm ; marge gauche 8.0 mm
\pgc{-12.700mm}{0.000mm}{
\l{\cel{6}92}
\l{}
\l{\sou{cinika}. I. Qua apartenas a la skolo filozofiala quan fondis}
\l{\cel{3}Antistenes, tale nomizita pro ke olu esis establisita}
\l{\cel{3}komence en Cinosargos, pre-urbo di la urbo Athinai. -}
\l{\cel{3}II. Qua su lokizas super deco, la konveno-reguli, e pudo-}
\l{\cel{3}ro. - DEFIS.}
\l{}
\l{\sou{cinocefalo}. (\sou{zool.}) Genero de simio, di qua la muzelo esas}
\l{\cel{3}longa quale che la hundo. - L.\cel{1}\sou{cynocephalus}. - EFIS.}
\l{}
\l{\sou{cintilo}. Peceto flama, lumoza, qua su separas de korpo}
\l{\cel{3}kombustata. - EFIS.}
\l{}
\l{\sou{cipo}. (\sou{arkitekt.}) Mi-kolono sen kapitelo, qua simulas kolono}
\l{\cel{3}ruptita, quan la antiqui erektis kom monumento funerala,}
\l{\cel{3}quan li lokizis sur la voyi por indikar la direcioni, od}
\l{\cel{3}en la loki konsakrita da ta o ca eventaji memorinda. -}
\l{\cel{3}DEFIRS.}
\l{}
\l{\sou{cipreso}. (\sou{bot.)}\cel{2}Arboro rezinala, de la familio \textquotedbl{}koniferi\textquotedbl{},}
\l{\cel{3}quan lua foliaro obskura igas selektar por plantaceso}
\l{\cel{3}en la tombeyi, apud la tombi. - L.\cel{1}\sou{cypressus}. - DEFIRS.}
\l{}
\l{\sou{ciprino}. (\sou{zool.)}\cel{2}Genero de fisho qua vivas en aquo sen-sala,}
\l{\cel{3}di qua la tipo esas la karpo, e di qua la koloro esas}
\l{\cel{3}pasable oranjea o redatra. - L.\cel{1}\sou{cyprinus}. - EFIS.}
\l{}
\l{\sou{cirajo}. Kompozajo nigra, quan onu extensas sur la ledro di}
\l{\cel{3}la pedovesti, di la harnesi, e quan onu frotas per bros-}
\l{\cel{3}ilo por igar lu brilanta. - FIS.}
\l{}
\l{\sou{circinata}. (\sou{dermatologio}). (Lezuri di pelo), grupigita cir-}
\l{\cel{3}klatre o cirklo-fragmentatre, di qua la centro esas sen-}
\l{\cel{3}domaja. - EFS.}
\l{}
\l{\sou{cirko}. I. Klozajo cirkla ube onu celebris (che la antiqui)}
\l{\cel{3}la ludi publika. - II. Sorto di teatro, cirkoforma, ube}
\l{\cel{3}la spektajo konsistas ye prodalaji di kavalko, di voltijo.}
\l{\cel{3}- DEFIRS.}
\l{}
\l{\sou{cirklo}. (\sou{geom.}) Parto de plano, di qua la limito esas kurvo}
\l{\cel{3}klozita di qua omna punti distas egale de punto interna}
\l{\cel{3}(centro). - EFIRS.}
\l{}
\l{\sou{cirkoncizar}.\cel{2}(\sou{trans}.)Distranchar la prepuco. - EFIS.}
\l{}
\l{\sou{cirkondar}. (\sou{trans.}) Esar cirkum. - II. Garnisar cirkume.}
\l{}
\l{\sou{cirkonflexo}. En la linguo Franca, Portugalana, e c. : signo}
\l{\cel{3}ortografiala( \textasciicircum{})\cel{2}quan onu lokizas super ula vokali longa}
\l{\cel{3}(\sur{\textasciicircum{}}{a}, \sur{\textasciicircum{}}{e}, \sur{\textasciicircum{}}{i}, \sur{\textasciicircum{}}{o}).\cel{2}- DEFIRS.}
}
